% Canonical LaTeX formulation: two-batch (REF vs LIVE) uplift monitoring
% isomorphic to batch FSDS, with AUUC ranking layer + uplift effect layer.
% Compile via uplift_fsds_two_batch_standalone.tex

\section{Two-batch uplift monitoring (FSDS-isomorphic formulation)}

\paragraph{Monitoring contract.}
Fix two disjoint finite samples representing \textbf{calibration} and \textbf{deployment}:
\[
  \mathcal{D}_{\mathrm{ref}} = \{(X_i,T_i,Y_i)\}_{i\in\mathcal{I}_{\mathrm{ref}}},
  \qquad
  \mathcal{D}_{\mathrm{live}} = \{(X_j,T_j,Y_j)\}_{j\in\mathcal{I}_{\mathrm{live}}}.
\]
Split $\mathcal{D}_{\mathrm{ref}}$ into training and evaluation subsets
$\mathcal{D}_{\mathrm{ref}}^{\mathrm{tr}}\cup\mathcal{D}_{\mathrm{ref}}^{\mathrm{ev}}$ (held out for in-batch AUUC).
The \textbf{uplift ranker} $\hat\tau:\mathbb{R}^p\to\mathbb{R}$ is fit only on $\mathcal{D}_{\mathrm{ref}}^{\mathrm{tr}}$ and then \textbf{frozen} when scoring $\mathcal{D}_{\mathrm{live}}$.
No refit on LIVE is permitted in the monitoring layer except inside LOCO ablations (always retrain on REF train only).

\paragraph{Batch vs treatment.}
Introduce batch indicator $W\in\{0,1\}$ with $W=0$ on $\mathcal{D}_{\mathrm{ref}}$ and $W=1$ on $\mathcal{D}_{\mathrm{live}}$.
This is \textbf{not} the randomised treatment $T\in\{0,1\}$ used in the uplift experiment.
All FSDS-style batch tests use $(X,W)$ or stacked $(X,Y,W)$; uplift curves use $(Y,T,\hat\tau(X))$ within each batch.

\paragraph{Isomorphism to predictive batch FSDS.}
\begin{center}
\small
\begin{tabular}{@{}lll@{}}
\toprule
\textbf{Predictive FSDS} & \textbf{Two-batch uplift monitor} & \textbf{Object} \\
\midrule
Global loss / calibration on LIVE & $\mathrm{AUUC}_{\mathrm{live}}$, $\Delta_{\mathrm{AUUC}}$ & ranking quality \\
Global MMD$^2$ on $X$ & $\widehat{\mathrm{MMD}}^2(X_{\mathrm{ref}},X_{\mathrm{live}})$ & covariate shift \\
Domain classifier AUC$(X)$ & $\widehat{\mathrm{AUC}}_{\mathrm{dom}}^{\mathrm{RF}}(X;W)$ & same \\
Effective sample size (overlap) & $\mathrm{ESS}_{\mathrm{ovlp}}(\hat e(W\mid X))$ & comparable support \\
LOGO-MMD on feature groups & LOGO-MMD$(g)$ on $X$ & localised $X$ shift \\
Leave-one-group-out on outcome model & LOCO--AUUC$(g)$ & localised ranking drivers \\
Permute--refit on batch (PO-risk) & DRPerm on $(X,Y,W)$ & $Y\mid X$ shift \\
Post-hoc population subset & Domain-quintile AUUC on LIVE & user-slice ranking \\
\bottomrule
\end{tabular}
\end{center}

\subsection{Validity gates (batch geometry)}

\paragraph{Covariate shift on $X$.}
Let $\hat{\phi}$ be an RKHS mean embedding distance or unbiased MMD$^2$ estimator:
\[
  \widehat{\mathrm{MMD}}^2
  = \widehat{\mathrm{MMD}}^2\bigl(\{X_i\}_{i\in\mathcal{I}_{\mathrm{ref}}^{\mathrm{ev}}},
  \{X_j\}_{j\in\mathcal{I}_{\mathrm{live}}}\bigr).
\]
Train a domain classifier $\hat d_W:\mathbb{R}^p\to[0,1]$ on stacked labelled pairs $(X,W)$ and report
$\widehat{\mathrm{AUC}}_{\mathrm{dom}}^{\mathrm{RF}}=\widehat{\mathrm{AUC}}(\hat d_W(X),W)$.

\paragraph{Batch propensity overlap.}
On the stacked sample, estimate $\hat e(w\mid x)=\mathbb P(W=1\mid X=x)$ and define overlap weights
$v_i = 1/(\hat e_i(1-\hat e_i))$ with clipping.
The effective sample size fraction is
\[
  \mathrm{ESS}_{\mathrm{ovlp}}
  = \frac{\bigl(\sum_i v_i\bigr)^2}{\sum_i v_i^2}\cdot\frac{1}{n_{\mathrm{ref}}+n_{\mathrm{live}}}.
\]
Monitoring attributions (LOCO--AUUC, slice tables) require $\mathrm{ESS}_{\mathrm{ovlp}}$ above a policy threshold; otherwise emit \textsc{RefreshRef} before subset rules.

\paragraph{Randomisation check on $T$.}
Within each batch $w\in\{0,1\}$, test treatment balance on $T$ (e.g.\ Pearson $\chi^2$ against design ratio $\pi$).
Failure $\Rightarrow$ invalid uplift curves; halt AUUC interpretation.

\subsection{Layer 1: ranking monitor (AUUC)}

\paragraph{Uplift ranker.}
Let $\mathcal{A}$ be a meta-learner factory (T-/X-/R-learner or registry nuisance models).
Fit $\hat\tau$ on $\mathcal{D}_{\mathrm{ref}}^{\mathrm{tr}}$:
\[
  \hat\tau \leftarrow \mathcal{A}\bigl(\mathcal{D}_{\mathrm{ref}}^{\mathrm{tr}}\bigr).
\]
Scores on evaluation windows:
\[
  \hat\tau_i^{\mathrm{ref}} = \hat\tau(X_i),\ i\in\mathcal{I}_{\mathrm{ref}}^{\mathrm{ev}};
  \qquad
  \hat\tau_j^{\mathrm{live}} = \hat\tau(X_j),\ j\in\mathcal{I}_{\mathrm{live}}.
\]

\paragraph{Area under the uplift curve (AUUC).}
Sort users by $\hat\tau$ descending; let $n_t(k)$, $n_c(k)$ be treated/control counts among the top-$k$ fraction.
With cumulative treated/control response sums $S_t(k)$, $S_c(k)$, define the uplift curve
$U(k)=\frac{S_t(k)}{n_t(k)}-\frac{S_c(k)}{n_c(k)}$ (with standard sklift conventions on ties and empty strata).
Then
\[
  \mathrm{AUUC}_{\mathrm{ref}}
  = \int_0^1 U_{\mathrm{ref}}(k)\,dk
  \quad\text{on}\quad \mathcal{D}_{\mathrm{ref}}^{\mathrm{ev}},
\]
and analogously $\mathrm{AUUC}_{\mathrm{live}}$ on $\mathcal{D}_{\mathrm{live}}$.
The \textbf{two-batch ranking alarm} is
\[
  \Delta_{\mathrm{AUUC}} = \mathrm{AUUC}_{\mathrm{ref}} - \mathrm{AUUC}_{\mathrm{live}}.
\]
Large $|\Delta_{\mathrm{AUUC}}|$ under stable gates indicates deployment ranking degradation relative to REF calibration.

\paragraph{LOCO--AUUC (feature-localised ranking).}
Let $\mathcal{G}=\{g_1,\ldots,g_G\}$ be optional feature groups; if $\mathcal{G}=\emptyset$, each single coordinate defines a group.
For each $g\in\mathcal{G}$, let $X^{(-g)}$ denote features with group $g$ removed.
Retrain $\hat\tau^{(-g)} \leftarrow \mathcal{A}(\mathcal{D}_{\mathrm{ref}}^{\mathrm{tr}}, X^{(-g)})$ and recompute AUUC on both windows:
\[
  \mathrm{drop}_{\mathrm{ref}}^{(g)}
  = \mathrm{AUUC}_{\mathrm{ref}}^{\mathrm{full}} - \mathrm{AUUC}_{\mathrm{ref}}^{(-g)},\quad
  \mathrm{drop}_{\mathrm{live}}^{(g)}
  = \mathrm{AUUC}_{\mathrm{live}}^{\mathrm{full}} - \mathrm{AUUC}_{\mathrm{live}}^{(-g)}.
\]
Define $\mathrm{drop\_gap}^{(g)}=\mathrm{drop}_{\mathrm{live}}^{(g)}-\mathrm{drop}_{\mathrm{ref}}^{(g)}$.
Large $|\mathrm{drop}_{\mathrm{live}}^{(g)}|$ localises LIVE ranking sensitivity to group $g$;
the vector $(\mathrm{drop}_{\mathrm{ref}}^{(g)})_g$ vs $(\mathrm{drop}_{\mathrm{live}}^{(g)})_g$ compared by Spearman $\rho$ distinguishes covariate-driven vs concept-driven reallocation of feature importance.

\paragraph{User slices on LIVE (domain quintiles).}
Train batch propensity on stacked covariates:
\[
  \hat s(x) = \mathbb P(W=1\mid X=x)
  \quad\text{fit on}\quad
  \{(X_i,W_i)\}_{i\in\mathcal{I}_{\mathrm{ref}}^{\mathrm{pool}}\cup\mathcal{I}_{\mathrm{live}}},
\]
where $\mathcal{I}_{\mathrm{ref}}^{\mathrm{pool}}$ is a REF pool (eval $+$ optional train cap).
On LIVE only, form quantile bins $Q_1,\ldots,Q_K$ by $\hat s(X_j)$ and compute
\[
  \mathrm{AUUC}_k = \mathrm{AUUC}\bigl(\mathcal{D}_{\mathrm{live}}\cap Q_k\bigr),\qquad k=1,\ldots,K.
\]
For slices $A,B\subseteq\mathcal{I}_{\mathrm{live}}$, pairwise ranking contrast
$\Delta_{\mathrm{AUUC}}^{A,B}=\mathrm{AUUC}_A-\mathrm{AUUC}_B$; policy ranks pairs by $|\Delta_{\mathrm{AUUC}}^{A,B}|$ for cap/continue rules.

\paragraph{Overlap-support slice.}
Let $\mathcal{L}_{\mathrm{ovlp}}=\{j\in\mathcal{I}_{\mathrm{live}}:\ \hat e(W=1\mid X_j)\in[\ell,u]\}$.
Report $\mathrm{AUUC}_{\mathrm{ovlp}}$ on $\mathcal{L}_{\mathrm{ovlp}}$.
If $\mathrm{AUUC}_{\mathrm{ovlp}}\gg \mathrm{AUUC}_{\mathrm{live}}$, global $\Delta_{\mathrm{AUUC}}$ is partially explained by batch mix off the shared support.

\subsection{Layer 2: uplift level on LIVE slices (effect layer)}

\paragraph{Observed slice uplift.}
On each quintile $Q_k\subseteq\mathcal{D}_{\mathrm{live}}$ with sufficient $T$-support,
\[
  \widehat{\tau}^{\mathrm{obs}}_k
  = \bar Y_{T=1,k}-\bar Y_{T=0,k}.
\]

\paragraph{Model slice uplift level.}
Using the same frozen REF ranker,
\[
  \bar{\hat\tau}_k = \frac{1}{|Q_k|}\sum_{j\in Q_k} \hat\tau(X_j).
\]

\paragraph{Pairwise effect contrasts (LIVE only).}
For quintiles $Q_a,Q_b$,
\[
  \Delta\widehat{\tau}^{\mathrm{obs}}_{a,b}
  = \widehat{\tau}^{\mathrm{obs}}_a - \widehat{\tau}^{\mathrm{obs}}_b,
  \qquad
  \Delta\bar{\hat\tau}_{a,b} = \bar{\hat\tau}_a - \bar{\hat\tau}_b.
\]
Layer~2 policy uses $|\Delta\widehat{\tau}^{\mathrm{obs}}_{a,b}|$ to prioritise budget reallocation \emph{conditional on concept rejection} (below).
Alignment of $\bar{\hat\tau}_k$ with $\widehat{\tau}^{\mathrm{obs}}_k$ across $k$ informs \textsc{Relearn} vs \textsc{Reallocate}.

\paragraph{Layer separation (do not conflate).}
Layer~1 measures \emph{sorting quality} of $\hat\tau$ on LIVE; Layer~2 measures \emph{level} of treatment contrast on LIVE slices.
A slice may exhibit high $\widehat{\tau}^{\mathrm{obs}}_k$ and low $\mathrm{AUUC}_k$ simultaneously; caps follow Layer~1, budget shifts on Layer~2 require stronger concept evidence.

\subsection{Concept axis: PO-risk on stacked batches}

\paragraph{Pseudo-outcome risk.}
Stack $\mathcal{D}=\mathcal{D}_{\mathrm{ref}}\cup\mathcal{D}_{\mathrm{live}}$ with batch labels $W$.
Cross-fit nuisances $\hat\mu(X)$, $\hat e(W\mid X)$; form pseudo-outcomes $\hat\eta=(Y-\hat\mu)(W-\hat e)$ and fit $\hat\tau_{\mathrm{po}}(X)$ by regressing $\hat\eta$ on $X$.
The observed PO-risk statistic is $R=\mathbb E[\hat\tau_{\mathrm{po}}(X)^2]$ (sample mean on $\mathcal{D}$).

\paragraph{Permute--refit test (DRPerm).}
For $b=1,\ldots,B$, permute batch labels $W^{(b)}$, refit nuisances and $\hat\tau_{\mathrm{po}}^{(b)}$, store $R^{(b)}$.
Two-sided Monte Carlo $p$-value:
\[
  p = \frac{1+\sum_{b=1}^B \mathbf 1\{ R^{(b)}\ge R\}}{B+1}.
\]
Reject at level $\alpha$ $\Rightarrow$ interpret large $\Delta_{\mathrm{AUUC}}$ with flat $\widehat{\mathrm{MMD}}^2$ as $Y\mid X$ (and treatment interaction) shift; enable Layer~2 REALLOCATE rules on empirical ATE pairs.

\subsection{Regime classification (two axes)}

\begin{center}
\small
\begin{tabular}{@{}p{0.22\linewidth}cc@{}}
\toprule
 & $Y\mid X$ stable ($p>\alpha$) & $Y\mid X$ shifted ($p\le\alpha$) \\
\midrule
$X$ stable (MMD, AUC flat) &
Layer~1: LOCO + quintile AUUC &
Layer~1 + Layer~2 ATE; PO-LOCO optional \\
\midrule
$X$ shifted (MMD or AUC $\uparrow$) &
LOGO-MMD$(g)$, quintile AUUC, overlap &
Stratify REF; gates then LOCO + L2 \\
\bottomrule
\end{tabular}
\end{center}

If $\mathrm{ESS}_{\mathrm{ovlp}}$ is low, suspend LOCO and slice effect rules; emit refresh/match REF.

\subsection{Monitoring map (formal pipeline)}

\begin{enumerate}
  \item \textbf{Gates:} SRM$(T)$ per batch; $\widehat{\mathrm{MMD}}^2$, $\widehat{\mathrm{AUC}}_{\mathrm{dom}}$; $\mathrm{ESS}_{\mathrm{ovlp}}$.
  \item \textbf{Layer~1 global:} $\mathrm{AUUC}_{\mathrm{ref}}$, $\mathrm{AUUC}_{\mathrm{live}}$, $\Delta_{\mathrm{AUUC}}$.
  \item \textbf{Layer~1 local:} LOCO--AUUC$(g)$; $\mathrm{AUUC}_k$; pairwise $|\Delta_{\mathrm{AUUC}}^{A,B}|$; optional $\mathrm{AUUC}_{\mathrm{ovlp}}$.
  \item \textbf{Concept:} DRPerm $p$ on stacked $(X,Y,W)$.
  \item \textbf{Layer~2 local:} $\widehat{\tau}^{\mathrm{obs}}_k$, $\bar{\hat\tau}_k$; pairwise $|\Delta\widehat{\tau}^{\mathrm{obs}}_{a,b}|$.
  \item \textbf{Policy:} map statistics to $\{\textsc{Reallocate},\textsc{HoldFeature},\textsc{RefreshRef},\textsc{Relearn},\textsc{Monitor}\}$.
\end{enumerate}

\paragraph{Policy mapping (examples).}
Let $\tau_{\mathrm{AUUC}}$ be AUUC SLA, $\delta_{\mathrm{mmd}}$ MMD threshold, $\epsilon_{\mathrm{ess}}$ ESS floor.
\begin{align*}
  \mathrm{ESS}_{\mathrm{ovlp}}<\epsilon_{\mathrm{ess}}
  &\Rightarrow \textsc{RefreshRef}.\\
  \mathrm{AUUC}_{\mathrm{live}}<\tau_{\mathrm{AUUC}}
  &\Rightarrow \textsc{Reallocate}\ \text{(global pause/cap)}.\\
  k^\star=\arg\min_k \mathrm{AUUC}_k
  &\Rightarrow \textsc{Reallocate}\ \text{(cap quintile $Q_{k^\star}$)}.\\
  \arg\max_{a,b}|\Delta\widehat{\tau}^{\mathrm{obs}}_{a,b}|,\ p\le\alpha
  &\Rightarrow \textsc{Reallocate}\ \text{(budget toward higher $\widehat{\tau}^{\mathrm{obs}}$ slice)}.\\
  \widehat{\mathrm{MMD}}^2\ge\delta_{\mathrm{mmd}},\ \mathrm{ESS}_{\mathrm{ovlp}}\ \text{OK}
  &\Rightarrow \textsc{HoldFeature}\ \text{(LOGO-MMD / LOCO top group)}.\\
  p\le\alpha,\ \text{systematic } \bar{\hat\tau}_k\not\approx\widehat{\tau}^{\mathrm{obs}}_k
  &\Rightarrow \textsc{Relearn}\ \text{after REALLOCATE trial}.
\end{align*}

\subsection{Two-batch information flow}

\begin{center}
\begin{tabular}{@{}c|c|c@{}}
\toprule
 & REF & LIVE \\
\midrule
Fit $\hat\tau$ & $\mathcal{D}_{\mathrm{ref}}^{\mathrm{tr}}$ & --- \\
AUUC eval & $\mathcal{D}_{\mathrm{ref}}^{\mathrm{ev}}$ & $\mathcal{D}_{\mathrm{live}}$ \\
LOCO retrain & $\mathcal{D}_{\mathrm{ref}}^{\mathrm{tr}}$ (per drop) & score only \\
Train $\hat s,\hat e$ for slices & $\mathcal{I}_{\mathrm{ref}}^{\mathrm{pool}}\cup\mathcal{I}_{\mathrm{live}}$ & slice LIVE rows \\
Empirical ATE & --- & $Q_k\subseteq\mathcal{D}_{\mathrm{live}}$ \\
PO-risk / DRPerm & \multicolumn{2}{c}{stacked $\mathcal{D}_{\mathrm{ref}}\cup\mathcal{D}_{\mathrm{live}}$} \\
\bottomrule
\end{tabular}
\end{center}

\paragraph{Implementation.}
\texttt{loco\_auuc.loco\_auuc\_monitor} (Layer~1 LOCO);
\texttt{run\_uplift\_subset\_localization} (gates, quintiles, pairwise, rules);
\texttt{run\_uplift\_subset\_benchmark.py} (four REF/LIVE scenarios).
Narrative playbook: \verb|docs/FSDS_UPLIFT_TWO_BATCH_LANDING.md|.
